% title: 공간벡터와 좌표축
% description: 3차원 좌표축 x, y, z와 벡터 v = (v_x, v_y, v_z). xy평면으로 내린 점선 투영과 성분 표시
% packages:
% libraries: arrows.meta
% 성분 \vx, \vy, \vz만 바꾼다. 비스듬한 투영: x축은 왼쪽 아래, y축은 오른쪽, z축은 위(tikzpicture 옵션의 x, y, z).
\begin{tikzpicture}[x={(-0.6cm,-0.45cm)}, y={(1cm,0cm)}, z={(0cm,1cm)}, scale=1.3,
                    >={Stealth[length=2.5mm]}]
  \def\vx{2}
  \def\vy{3}
  \def\vz{2.5}
  \coordinate (O) at (0,0,0);
  \coordinate (P) at (\vx,\vy,\vz);
  \coordinate (Q) at (\vx,\vy,0);
  \draw[->] (O) -- (\vx+1.5,0,0) node[below left] {$x$};
  \draw[->] (O) -- (0,\vy+1.2,0) node[right] {$y$};
  \draw[->] (O) -- (0,0,\vz+1) node[above] {$z$};
  \draw[dashed, gray] (\vx,0,0) -- (Q) -- (0,\vy,0);
  \draw[dashed, gray] (O) -- (Q);
  \draw[dashed, gray] (Q) -- node[right] {$v_z$} (P);
  \draw[->, very thick, blue] (O) -- (P) node[above right] {$\vec{v}=(v_x,\,v_y,\,v_z)$};
  \fill (Q) circle[radius=1.2pt];
  \node[above left] at (\vx,0,0) {$v_x$};
  \node[above] at (0,\vy,0) {$v_y$};
  \node[left] at (O) {$O$};
\end{tikzpicture}
